\documentclass[11pt]{article}
\usepackage[T1]{fontenc}
\usepackage{lmodern}
\usepackage{amsmath,amssymb,amsthm}
\usepackage{booktabs}
\usepackage{enumitem}
\usepackage{graphicx}
\usepackage[margin=0.92in]{geometry}
\usepackage[colorlinks=true,linkcolor=blue,citecolor=blue,urlcolor=blue]{hyperref}

\newtheorem{theorem}{Theorem}
\newtheorem{lemma}{Lemma}
\newtheorem{proposition}{Proposition}
\newtheorem{remark}{Remark}
\newcommand{\R}{\mathbb R}
\newcommand{\E}{\mathbb E}
\newcommand{\Pp}{\mathbb P}

\title{A Nuclear-Laplace Reduction for the General\\
Smallest-Singular-Value Problem}
\author{Run \texttt{fa\_banach\_001}, agent \texttt{agent\_lane\_14}}
\date{19 July 2026}

\begin{document}
\maketitle

\section*{Status}

Conditional counterexample reduction, likely valid.  The matrix construction,
the exact Pfaffian formula, and the implication to a counterexample are proved.
One scalar asymptotic dependency remains: an explicit Pfaffian ratio must be
shown to be unbounded.  Exact and high-precision computations support that
dependency but are not a proof.

\section*{Source question}

Livshyts, Fernandez, and Mui prove
\[
 \Pp\{s_n(A)\leq \varepsilon/\sqrt n\}
 \leq C\varepsilon+e^{-cn}
\]
for square unconditional isotropic log-concave matrices, and ask whether the
same behavior persists for general square isotropic log-concave matrix
laws~\cite{LFM}.  Their discussion on PDF page~4 is reproduced below.

\begin{center}
\includegraphics[width=0.96\linewidth]{figures/open_problem_crop.png}
\end{center}

\section*{Idea of the proof}

The nuclear norm is the convex matrix analogue of the $\ell_1$ norm, so the
law with density $e^{-\|A\|_*}$ is an especially natural place for excessive
mass near low rank.  Orthogonal invariance makes its covariance scalar, hence
isotropic normalization costs only one scalar.

For real square matrices the ordered singular values have density
\[
 e^{-\sum x_i}\prod_{i<j}(x_j^2-x_i^2).
\]
The density of the minimum at zero is therefore a ratio $d_n=B_n/Z_n$ of two
partition functions.  De Bruijn's identity turns both into completely
explicit Pfaffians.  More importantly, shifting all remaining singular values
by $x$ proves the quantitative inequality $g_n(x)\geq e^{-nx}d_n$.  Thus an
unbounded boundary ratio creates probability at least a diverging multiple of
$1/n$ on a window whose isotropic parameter is at most $1/n$.  This is
incompatible with any universal linear small-ball constant.

\section*{The ensemble and the conditional result}

Let $X_n$ be a random real $n\times n$ matrix with density
\begin{equation}\label{eq:matrix-law}
 d\mu_n(A)=\mathcal Z_n^{-1}e^{-\|A\|_*}\,dA.
\end{equation}
The density is log-concave and invariant under $A\mapsto UAV$ for
$U,V\in O(n)$.  It is centered, and its entry covariance is $v_n I_{n^2}$ for
some $v_n>0$.  Put $\widetilde X_n=X_n/\sqrt{v_n}$.

For $0<x_1<\cdots<x_n$, define
\begin{align}
 Z_n&=\int e^{-\sum_{i=1}^n x_i}
       \prod_{1\leq i<j\leq n}(x_j^2-x_i^2)\,dx_1\cdots dx_n,\label{eq:Z}\\
 B_n&=\int_{0<y_1<\cdots<y_{n-1}} e^{-\sum_{j=1}^{n-1}y_j}
       \prod_{j=1}^{n-1}y_j^2
       \prod_{1\leq i<j\leq n-1}(y_j^2-y_i^2)\,dy.\label{eq:B}
\end{align}
For $n=1$, the empty integral defining $B_1$ is $1$.  Write
$d_n=B_n/Z_n$.

\begin{theorem}[conditional counterexample]\label{thm:conditional}
If $d_n$ is unbounded along a sequence of dimensions, then no absolute
constants $C,c>0$ can satisfy
\[
 \Pp\{s_n(A)\leq\varepsilon/\sqrt n\}
 \leq C\varepsilon+e^{-cn}
\]
for every square isotropic log-concave random matrix $A$ and every
$\varepsilon>0$.  The isotropic matrices $\widetilde X_n$ from
\eqref{eq:matrix-law} give the counterexample sequence.
\end{theorem}

\section*{Proof of the conditional theorem}

The real square SVD change of variables shows that the ordered singular values
of $X_n$ have joint density $Z_n^{-1}$ times the integrand in
\eqref{eq:Z}; the orthogonal-group constant cancels in the normalization.
Let $g_n$ be the density of the least singular value.  Directly,
\begin{equation}\label{eq:g}
 g_n(x)=\frac{e^{-x}}{Z_n}
 \int_{x<y_1<\cdots<y_{n-1}}e^{-\sum y_j}
 \prod_j(y_j^2-x^2)\prod_{i<j}(y_j^2-y_i^2)\,dy.
\end{equation}
In particular, $g_n(0)=B_n/Z_n=d_n$.

\begin{lemma}[a quantitative boundary window]\label{lem:window}
For every $x\geq0$,
\[
 g_n(x)\geq e^{-nx}d_n.
\]
\end{lemma}
\begin{proof}
In \eqref{eq:g}, substitute $y_j=z_j+x$.  The exponential contributes
$e^{-nx}$.  For $i<j$,
\[
 (z_j+x)^2-(z_i+x)^2=(z_j-z_i)(z_j+z_i+2x)
 \geq z_j^2-z_i^2,
\]
and
\[
 (z_j+x)^2-x^2=z_j(z_j+2x)\geq z_j^2.
\]
The shifted integral is therefore at least $B_n$.
\end{proof}

It remains to track isotropic normalization.  Orthogonal row and column sign
flips kill every covariance between distinct entries, while row and column
permutations make all entry variances equal.  Thus
$\operatorname{Cov}(\operatorname{vec}X_n)=v_nI$.

There is also a useful dimension-free lower bound on $v_n$.  Radial
integration for any norm in dimension $D=n^2$ gives
\[
 \E\|X_n\|_*^2=D(D+1)=n^2(n^2+1).
\]
Since $\|A\|_*\leq\sqrt n\|A\|_F$,
\[
 n^2v_n=\E\|X_n\|_F^2
 \geq \frac1n\E\|X_n\|_*^2,
 \qquad v_n\geq n+\frac1n.\label{eq:variance}
\]

Integrating Lemma~\ref{lem:window} on $[0,1/n]$ yields
\begin{equation}\label{eq:lowerprob}
 \Pp\{s_n(X_n)\leq1/n\}
 \geq \frac{1-e^{-1}}{n}d_n.
\end{equation}
For
\[
 \varepsilon_n=\frac{1}{\sqrt n\sqrt{v_n}},
\]
the event in \eqref{eq:lowerprob} is exactly
$\{s_n(\widetilde X_n)\leq\varepsilon_n/\sqrt n\}$, and
the preceding variance estimate gives $\varepsilon_n\leq1/n$.  If universal $C,c$
existed, \eqref{eq:lowerprob} would imply
\[
 (1-e^{-1})d_n\leq C+n e^{-cn}.
\]
This contradicts unboundedness of $d_n$ and proves
Theorem~\ref{thm:conditional}.  \qed

\section*{Exact Pfaffian reduction}

The remaining dependency is not a high-dimensional integral.  It is an
explicit Pfaffian ratio.  For integers $0\leq i<j$, set
\begin{align}
 q_{ij}
 &=\frac{1}{(2i)!(2j)!}
 \int_{0<x<y}e^{-x-y}
   \bigl(x^{2i}y^{2j}-x^{2j}y^{2i}\bigr)\,dx\,dy\label{eq:qint}\\
 &=2I_{1/2}(2i+1,2j+1)-1,\label{eq:qbeta}
\end{align}
and put $q_{ji}=-q_{ij}$, $q_{ii}=0$.  For an ordered finite index set
$S$, let $\operatorname{DPf}(S)$ denote the Pfaffian of $(q_{ij})_{i,j\in S}$
when $|S|$ is even.  When $|S|$ is odd, augment this skew matrix by a last
column of ones and a last row of minus ones before taking the Pfaffian.

\begin{proposition}[de Bruijn formula]\label{prop:pf}
With the sign convention in \eqref{eq:qint},
\[
 d_n=\frac{\operatorname{DPf}(\{1,\ldots,n-1\})}
           {\operatorname{DPf}(\{0,\ldots,n-1\})}.
\]
\end{proposition}
\begin{proof}
Write the Vandermonde factor as
\[
 \prod_{i<j}(x_j^2-x_i^2)=
 \det[x_j^{2(i-1)}]_{i,j=1}^n.
\]
De Bruijn's ordered-chamber integration identity expresses its integral
against $\prod e^{-x_j}dx_j$ as the Pfaffian of the skew moments in
\eqref{eq:qint}; for odd size the unpaired moments are
$\int_0^\infty x^{2i}e^{-x}dx=(2i)!$, producing the one-vector after
normalization.  The factors $\prod_{i=0}^{n-1}(2i)!$ occur in $Z_n$.
For $B_n$, the basis exponents are $2,4,\ldots,2(n-1)$, so the factors are
$\prod_{i=1}^{n-1}(2i)!$.  Their ratio is $1/0!=1$, proving the formula.
Finally, \eqref{eq:qbeta} follows by interpreting normalized gamma integrals:
if $G_i,G_j$ are independent gamma variables of shapes $2i+1,2j+1$, then
$G_i/(G_i+G_j)$ is beta distributed.
\end{proof}

\section*{Conditional dependencies and computation}

The only unproved statement is
\begin{equation}\label{eq:dependency}
 \sup_n d_n=\infty.
\end{equation}
No rate is needed.  The supplied exact-fraction program implements
Proposition~\ref{prop:pf} and gives:
\begin{center}
\begin{tabular}{@{}rrrrrr@{}}
\toprule
$n$ & 1 & 10 & 40 & 100 & 200\\
\midrule
$d_n$ & 1.000000 & 2.038100 & 2.577737 & 2.917386 & 3.168924\\
\bottomrule
\end{tabular}
\end{center}
Values through $n=40$ were computed exactly as rational numbers; a separate
high-precision implementation continued the audit through $n=200$.  The data
are consistent with growth comparable to $\sqrt{\log n}$, but neither that fit
nor the weaker dependency \eqref{eq:dependency} is claimed as proved.

\section*{Limitations and novelty check}

This packet is conditional, not a claimed disproof.  A bounded search used the
exact source title and arXiv id together with ``smallest singular value
Schatten 1 ball'', ``nuclear norm random matrix hard edge'', and close
Pfaffian/partition-function phrases.  The closest inspected sources concerned
macroscopic Schatten-ball equilibrium measures~\cite{DFGZ,KPT}, the critical
Freud ensemble~\cite{CKM}, and hard-edge universality for other Laguerre-type
weights~\cite{DGKV}.  No source was found that states \eqref{eq:dependency} or
this counterexample reduction.

\section*{Human-review recommendation}

The matrix argument should be checked by a specialist in asymptotic convex
geometry; the remaining ratio should be sent to a specialist in Pfaffian
ensembles or skew-orthogonal polynomials.  Proving any divergence of the ratio
would immediately promote this packet to a full counterexample.

\begin{thebibliography}{9}
\bibitem{LFM}
G. V. Livshyts, M. Fernandez V, and S. Mui,
\emph{On the Smallest Singular Value of Log-Concave Random Matrices},
arXiv:2508.17745 (2025).

\bibitem{DFGZ}
B. Dadoun, M. Fradelizi, O. Gu\'edon, and P.-A. Zitt,
\emph{Asymptotics of the inertia moments and the variance conjecture in
Schatten balls}, J. Funct. Anal. 284 (2023), 109752; arXiv:2111.07803.

\bibitem{KPT}
Z. Kabluchko, J. Prochno, and C. Th\"ale,
\emph{Exact asymptotic volume and volume ratio of Schatten unit balls},
J. Approx. Theory 257 (2020), 105457; arXiv:1804.03467.

\bibitem{CKM}
T. Claeys, I. Krasovsky, and O. Minakov,
\emph{Weak and strong confinement in the Freud random matrix ensemble and
gap probabilities}, Comm. Math. Phys. 402 (2023), 285--340;
arXiv:2209.07253.

\bibitem{DGKV}
P. Deift, D. Gioev, T. Kriecherbauer, and M. Vanlessen,
\emph{Universality for orthogonal and symplectic Laguerre-type ensembles},
J. Stat. Phys. 129 (2007), 949--1053; arXiv:math-ph/0612007.
\end{thebibliography}

\end{document}
